% !TEX root = ../main.tex
\section{Versuchsaufbau und Messung}
\label{sec:aufbau}

\subsection{Leistungskreis}

Der Aufbau war fertig verkabelt (Abbildung~\ref{fig:blockschaltbild}). Die Asynchronmaschine (ASM) wird aus dem Netz~1 gespeist und treibt über die gemeinsame Welle die Gleichstrommaschine (GM). Der Anker der GM liegt am Zwischenkreis des einphasigen Wechselrichters (WR). Die Zwischenkreisspannung $\Udc$ wird über den Erregerstrom der GM (Klemmen F1/F2) eingestellt, die Drehzahl wird von der ASM praktisch konstant gehalten. Der Wechselrichterausgang ist über die Kopplungsinduktivität und einen Trenntransformator mit dem Netz~2 verbunden.

\begin{figure}[htbp]
  \centering
  \begin{tikzpicture}[
      blk/.style={draw,rectangle,minimum height=1.1cm,minimum width=1.4cm,align=center,font=\footnotesize},
      >={Stealth[length=2mm]},node distance=0.8cm]
    \node[blk] (n1) {Netz 1\\ $3\times\SI{400}{\volt}$};
    \node[blk,right=of n1] (asm) {ASM\\ \SI{5.5}{\kilo\watt}};
    \node[blk,right=of asm] (gm) {GM\\ \SI{6}{\kilo\watt}};
    \node[blk,right=of gm] (wr) {WR\\ IGBT-\\ Vollbrücke};
    \node[blk,right=of wr] (l) {$L$\\ \SI{70}{\milli\henry}};
    \node[blk,right=of l] (tr) {Trafo\\ $400/230$\,V};
    \node[blk,right=of tr] (n2) {Netz 2\\ \SI{230}{\volt}};
    \node[blk,below=0.8cm of gm] (err) {Erregung\\ F1/F2};
    \draw (n1) -- (asm);
    \draw[double distance=1.5pt] (asm) -- (gm);
    \draw (gm) -- (wr);
    \draw (wr) -- (l);
    \draw (l) -- (tr);
    \draw (tr) -- (n2);
    \draw[->] (err) -- (gm);
    \node[font=\scriptsize,above] at ($(gm.east)!0.5!(wr.west)$) {$\Udc$};
    \node[font=\scriptsize,above] at ($(wr.east)!0.5!(l.west)$) {$\uWR$};
    \node[font=\scriptsize,above] at ($(tr.east)!0.5!(n2.west)$) {$\uN$};
  \end{tikzpicture}
  \caption{Blockschaltbild des Versuchsaufbaus. Die Welle zwischen ASM und GM ist doppelt dargestellt; $\Udc$ wird über den Erregerstrom der GM eingestellt.}
  \label{fig:blockschaltbild}
\end{figure}

Der Wechselrichter ist ein synchronisierbarer Einphasenwechselrichter mit IGBT-Vollbrücke und antiparallelen Dioden, geteiltem Zwischenkreiskondensator (Mittelpunkt M.P.) und einem antiparallelen Koppelschalter auf der AC-Seite. Er hängt auf der \SI{230}{\volt}-Seite des Trenntransformators; bei einer Zwischenkreisspannung von $\Udc=\SI{259}{\volt}$ wurde eine Grundschwingung der Wechselrichterspannung von $U_1\approx\SI{233}{\volt}$ gemessen (Abschnitt~\ref{sec:messung611}).

Am Bedienfeld lassen sich der Winkel $\Th$ zwischen Wechselrichter- und Netzspannung (Tasten $\leftarrow$ $\rightarrow$), die Gleichspannung (Tasten $\uparrow$ $\downarrow$, wirkt auf die GM-Erregung), das Impulsmuster, die Synchronisation und das Koppeln bzw.\ Entkoppeln bedienen. Zur Verfügung stehen die geschlossene Rechteckschwingung (Grundfrequenztaktung), die Muster mit eliminierter 3.\ bis 5., 3.\ bis 9.\ und 3.\ bis 21.\ Oberschwingung sowie das Trägerverfahren (sinusbewertet) fein und grob.

\subsection{Geräte und Messpunkte}

Die Entkopplungsinduktivität hat \SI{70}{\milli\henry} (\SI{30}{\ampere}); bei \SI{50}{\hertz} ist $X=\omega L\approx\SI{22}{\ohm}$. Hinzu kommen ein Einphasentransformator $400/230$\,V (\SI{25}{\ampere}), eine Gleichstrommaschine (\SI{6}{\kilo\watt}, \SI{500}{\volt}) und eine Asynchronmaschine (\SI{5.5}{\kilo\watt}, $3\times\SI{400}{\volt}$). Gemessen wurde mit dem Voltech-Power-Analyzer PM\,3000 und seinem Spannungswahlschalter (Stellung «Netz»: $\uN$, Stellung «WR»: $\uWR$; der Strom ist immer $\iN$), dem Vierkanal-Oszilloskop Tektronix TPS2014 (isolierte Kanäle) mit Stromzange und Trennverstärkern, dem Multimeter Metra Hit 18S ($\Udc$) und einem optischen Drehzahlmesser. Tabelle~\ref{tab:messpunkte} nennt die Messpunkte.

\begin{table}[htbp]
  \centering
  \caption{Messpunkte der Grössen.}
  \label{tab:messpunkte}
  \small
  \begin{tabular}{@{}ll@{}}
    \toprule
    Grösse & Messort \\
    \midrule
    $\uN$ & Buchsen «Trenntransformator Sekundär» \\
    $\uWR$ & Buchsen direkt links der H-Brücke \\
    $\iN$ & Stromzange in der AC-Leitung \\
    $\Udc$ & Wechselrichter-Bedienfeld, Klemmen $+$ und $-$ \\
    $u_{\mathrm{Halbleiter}}$ & $+$ gegen obere $\uWR$-Buchse (oberer Schalter, linker Zweig) \\
    \bottomrule
  \end{tabular}
\end{table}

Für die Leistungsbetrachtung stand der Spannungswahlschalter des PM\,3000 auf «Netz». Nur für die Spannungsoberschwingungen von $\uWR$ wurde auf «WR» umgeschaltet; die Strom-Oberschwingungen wurden in Stellung «Netz» abgelesen.

\subsection{Vorzeichenkonvention}
\label{sec:vorzeichen}

Positive Wirkleistung $P>0$ bedeutet Leistungsfluss aus dem Netz~2 über Trenntransformator, Kopplungsinduktivität und Wechselrichter in die Gleichstrommaschine. Der Wechselrichter arbeitet dabei als Gleichrichter, die GM als Motor und die ASM als Generator, der in das Netz~1 speist. Bei $P<0$ wird Leistung in das Netz~2 zurückgespeist (Wechselrichterbetrieb, GM als Generator, ASM als Motor). Der Winkel $\Th$ ist der Nacheilwinkel der Grundschwingung von $\uWR$ gegenüber $\uN$; ein positives $\Th$ gehört zu $P>0$. Die Blindleistung $Q$ des PM\,3000 ist als $\sqrt{S^2-P^2}$ definiert und daher immer positiv; sie enthält die Verzerrungsblindleistung.

Die Winkelanzeige am Wechselrichter war gegenüber der in der Vorlesung gelernten Darstellung vorzeichenverkehrt. Dadurch wurden die Lastpunkte (Abschnitte~\ref{sec:messung621} und \ref{sec:messung64}) mit $P>0$ gefahren, also im Gleichrichterbetrieb mit der GM als Motor und der ASM als Generator, und nicht im Wechselrichterbetrieb. Nach Einschätzung der Gruppe ändert das an den Aussagen des Berichts nichts.
